\section{Computational Verification}

All constructions and claims relying on explicit enumeration are
verified by explicit computation.

\subsection{Overview}

A computational pipeline was implemented to construct the graphs
$\Gsixty$, $\Gthirty$, and $\Gfifteen$, compute transport invariants,
and verify all structural and algebraic claims stated in this paper.

All computations are deterministic and based on explicit combinatorial
data derived from the dodecahedral flag structure.

\subsection{Graph construction and quotients}

The graph $\Gsixty$ is constructed from chamber states of the
dodecahedron. Two successive involutive identifications produce
$\Gthirty$ and $\Gfifteen$.

The following properties are verified:
\begin{itemize}
\item $\Gsixty$: $60$ vertices, $120$ edges, $4$-regular,
\item $\Gthirty$: $30$ vertices, $60$ edges, $4$-regular,
\item $\Gfifteen$: $15$ vertices, $30$ edges, $4$-regular,
\item all graphs are connected,
\item each quotient map is a $2$-fold covering.
\end{itemize}

\subsection{Identification of $\Gfifteen$}

The graph $\Gfifteen$ is verified to be isomorphic to the line graph of
the Petersen graph by matching vertex count, degree, and adjacency
structure.

\subsection{Transport invariant}

The function $\epsinv$ is computed by propagating signs along transport
paths from a fixed base chamber. Path independence is verified by
checking that all closed transport loops have trivial sign.

The induced edge sign function $\sigfun$ on $\Gfifteen$ is computed and
shown to have
\[
10 \text{ positive edges}, \qquad 20 \text{ negative edges}.
\]

\subsection{Cohomological classification}

All vertex switchings (with one vertex fixed to remove global symmetry)
are enumerated. The following are verified:
\begin{itemize}
\item no switching eliminates all negative edges,
\item the minimal support size is $6$,
\item exactly $40$ distinct minimal supports of size $6$ occur in the
gauge-fixed search.
\end{itemize}

\subsection{Sector incidence identity}

The sector--edge incidence matrix $M$ is constructed explicitly.
The identity
\[
MM^{\mathsf T} = A^3 + 2A^2 + 2I
\]
is verified by direct matrix computation.

\subsection{Reproducibility}

Details of the computational pipeline and reproducibility are given in
Appendices~D and~E.
